% !TEX root = ../algebra_02.tex

\section{Теорема Кронекера-Капелли}

\begin{Theorem}{Кронекера--Капелли}{}
	Система линейных уравнений $Ax = b$, где $A \in k^{m \times n}$, а $b \in k^m$ совместна тогда и только тогда, когда
	\[
	\mathrm{rk}(A \mid b) = \mathrm{rk}(A).
	\]
\end{Theorem}

\begin{proof}
	Распишем левую часть уравнения $Ax = b$:
	\[
	Ax = x_1 A[{:},1] + x_2 A[{:},2] + \cdots + x_n A[{:},n].
	\]
	Таким образом, система совместна (существуют такие числа $x_1, \dots, x_n$) $\Leftrightarrow$ столбец $b$ выражается в виде линейной комбинации столбцов матрицы $A$. Иными словами, вектор $b$ принадлежит линейной оболочке столбцов $A$:
	\[
	b \in \mathrm{Lin}(A[{:},1], \ldots, A[{:},n]).
	\]
	
	Это вложение равносильно тому, что добавление вектора $b$ к системе столбцов матрицы $A$ не увеличивает размерность их линейной оболочки:
	\[
	\dim \mathrm{Lin}(A[{:},1], \ldots, A[{:},n], b) = \dim \mathrm{Lin}(A[{:},1], \ldots, A[{:},n]).
	\]
	По определению столбцового ранга матрицы это означает, что ранг расширенной матрицы $(A \mid b)$ равен рангу исходной матрицы $A$, то есть $\mathrm{rk}(A \mid b) = \mathrm{rk}(A)$.
\end{proof}
